\section{Smooth-core folds and finite focus equations}
\label{app:harvesting-fold}

Threshold fitting uses a smooth-core cyclic fold with negative return
curvature and a nonsingular two-parameter fold equation. We derive these
inputs from finite focus equations and control the one-turn remainder
with the required state and parameter derivatives.

\subsection{The core chart and its focus coefficients}

At \(u=0\), use the equilibrium-preserving chart
\begin{equation}\label{eq:harvesting:core-chart}
 d=8\beta+\tfrac15+\nu,\qquad
 h=\tfrac{16}{5}\beta+\tfrac4{25}+\tfrac45\nu.
\end{equation}
The equilibrium is \((1/5,4/5)\), with
\[
 J=\begin{pmatrix}-1/5&-1/5\\16\beta&1/5+\nu\end{pmatrix},
 \qquad \operatorname{tr}J=\nu,\qquad
 \det J=\frac{80\beta-1-5\nu}{25}.
\]
The multiplicity-two weak focus of
\cite[Theorem~3.6]{huang2013} occurs at
\begin{equation}\label{eq:harvesting:bautin-point}
 \beta_0=\frac{\sqrt{51}-4}{100},\qquad
 d_0=\frac{-3+2\sqrt{51}}{25},\qquad
 h_0=\frac{4(1+\sqrt{51})}{125},\qquad \nu=0.
\end{equation}
The finite homogeneous equations below construct \(V=\mathcal R^2/2+\bigO(\mathcal R^3)\) in rotation
coordinates of radius \(\mathcal R\), and give
\begin{equation}\label{eq:harvesting:focus-data}
 k_4=\ell_1(\beta)
 =\frac{25\beta(2000\beta^2+160\beta-7)}{4(80\beta-1)},
 \qquad k_6\big|_{\beta=\beta_0}
 =-\frac{155}{48}-\frac{235\sqrt{51}}{576}<0.
\end{equation}
Thus \(\ell_1(\beta_0)=0\) and
\(\ell_1'(\beta_0)=17/2+4\sqrt{51}>0\). At the focus,
\(\dot V=k_6\mathcal R^6+\bigO(\mathcal R^7)\), so
\(\sqrt{2V}\) has nonzero negative quintic drift. The angular
velocity tends to
\(\omega_0=\sqrt{(-21+4\sqrt{51})/125}>0\); positive radial
scaling and angular clock therefore give \(\ell<0\) below.
The conjugate eigenvalues of \(J\) have real part \(\nu/2\), and
\((\nu,\beta)\mapsto(\nu/2,\ell_1(\beta))\) has rank two at
\((0,\beta_0)\).

At \(u=\nu=0\), put
\(\omega^2=(80\beta-1)/25\), \(X=x-1/5\), and
\(Z=(x-1/5+y-4/5)/(5\omega)\). Direct substitution gives
\begin{equation}\label{eq:harvesting:rotation}
 \dot X=-\omega Z+\Phi,\qquad
 \dot Z=\omega X+\Psi,\qquad \Phi=-5\omega XZ,
\end{equation}
where
\[
 \Psi=-\frac{
 80\omega X^2Z+(125\omega^2+5)X^2
 +(6\omega-250\omega^3)XZ
 +(125\omega^4+5\omega^2)Z^2}
 {16\omega(5X+1)}.
\]
Expand \((1+5X)^{-1}\) through the degrees needed below and denote
homogeneous terms by subscripts. For
\(\mathcal R^2=X^2+Z^2\), seek
\(V=\mathcal R^2/2+V_3+\cdots+V_6\). The equations are
\begin{equation}\label{eq:harvesting:homological}
 \mathcal L V_n=k_n\mathcal R^n-H_n,\qquad
 \mathcal L=\omega(-Z\partial_X+X\partial_Z),\qquad
 H_n=\sum_{j=2}^{n-1}
 \bigl(\partial_XV_j\,\Phi_{n-j+1}
       +\partial_ZV_j\,\Psi_{n-j+1}\bigr).
\end{equation}
Here \(V_2=\mathcal R^2/2\). On homogeneous polynomials of odd degree
\(n\), \(\mathcal L\) is invertible and \(k_n=0\); for even \(n\), take
the radial obstruction
\[
 k_n=\frac1{2\pi}\int_0^{2\pi}
 H_n(\cos\theta,\sin\theta)\,d\theta
\]
and set the \(X^n\) coefficient of \(V_n\) to zero to fix the radial
kernel. Solving \eqref{eq:harvesting:homological} successively in the
monomial basis for $n=3,4,5,6$ gives
$k_3=k_5=0$ and the coefficients~\eqref{eq:harvesting:focus-data}.
A program in Online Resource~2, also included with the source files,
reconstructs the rotated field and checks every coefficient residual of these
finite equations in exact arithmetic.

\subsection{Normal form and cyclic folds}

The focus coefficients supply the generalized Hopf unfolding. Its
weighted return turns the limiting fold into a periodic orbit with the
geometry required at the thresholds.

Take the finite normal form through degree
seven in a parameter-dependent complex coordinate. Its nonresonant
divisors stay bounded away from zero near the positive frequency
\(\omega_0\); hence the transformations and their required parameter
derivatives are smooth. The retained resonant degrees are one,
three, five and seven. Division by the positive angular speed and
normalization of one turn give smooth coefficients \(\eta_1,\eta_2\)
for the first two radial terms, with invertible
\((\nu,\beta)\mapsto(\eta_1,\eta_2)\) Jacobian. The fifth-degree
coefficient equals \(\ell\ne0\) at the focus. For the generalized Hopf normal
form see~\cite{kuznetsov2023}; its coefficients for ordinary and delay
equations are also derived by N.~A.~M.~Delmeire, M.~M.~Bosschaert and
Yu.~A.~Kuznetsov (preprint, 2025,
\href{https://arxiv.org/abs/2507.18469v1}{arXiv:2507.18469v1}).
Normalize the angular clock to one turn and put
\(r=\rho z\), \(\eta_1=\rho^4s_1\), \(\eta_2=\rho^2s_2\).
Then
\begin{equation}\label{eq:harvesting:weighted-flow}
 \frac{dz}{d\theta}
 =\rho^4(s_1z+s_2z^3+\ell z^5)+\bigO(\rho^6),\qquad \ell<0.
\end{equation}
On a fixed positive \(z\)-interval about one, the bound holds after
two \(z\) derivatives and one derivative in \((s_1,s_2)\), including
mixed derivatives. The quintic coefficient varies by \(\bigO(\rho^2)\),
giving \(\bigO(\rho^6)\) in the scaled equation, as does the retained
seventh-degree term; the smooth degree-eight field remainder gives
\(\bigO(\rho^7)\). Taylor remainders preserve these state bounds,
and scaled parameter differentiation introduces a factor \(\rho^4\)
or \(\rho^2\). The angular speed retains a positive uniform margin.
Integration over one turn and smooth finite-time flow dependence give
the smooth-core displacement
\begin{equation}\label{eq:harvesting:scaled-return}
 \rho^{-5}D(\rho z;\rho^4s_1,\rho^2s_2)
 =s_1z+s_2z^3+\ell z^5+\smallO(1)
\end{equation}
in the same norms, including angle-dependent remainders.
At \(z=1\), the limiting fold equations
\(s_1+s_2+\ell=0\), \(s_1+3s_2+5\ell=0\) have parameter
determinant two. Hence for every sufficiently small positive \(\rho\)
the implicit function theorem gives a cyclic fold with
\begin{equation}\label{eq:harvesting:fold-data}
 \eta_1=\ell\rho^4+\smallO(\rho^4),\qquad
 \eta_2=-2\ell\rho^2+\smallO(\rho^2),\qquad
 D_{rr}(\rho)=8\ell\rho^3+\smallO(\rho^3)\ne0.
\end{equation}
The \((D,D_r)\) parameter determinant is a nonzero
coordinate factor times \(2\rho^3+\smallO(\rho^3)\). Moreover,
\(\partial_\nu\eta_1\ne0\) at the generalized Hopf point and
\[
 \partial_\nu D(\rho)
 =\rho\partial_\nu\eta_1+\smallO(\rho)\ne0.
\]
These margins are strict at each fixed positive radius.
After division by \(\rho\), the fold orbit is \(C^2\)-close to a
Hopf ellipse. Its prey derivative is close to a nonzero sinusoid,
with strict sign away from two simple zeros and strict monotonicity
near them. Thus the orbit is embedded and nonstationary, has exactly
one nondegenerate prey minimum and maximum, and lies in a compact
subset of \(x,y>0\).
This is the fold used to fit the lower and upper thresholds in the
harvesting model.
